\documentclass[11pt,a4paper]{article}
\usepackage[utf8]{inputenc}
\usepackage[T1]{fontenc}
\usepackage[catalan]{babel}
\usepackage{lmodern}
\usepackage[sfdefault,lining]{FiraSans}
\usepackage[a4paper,top=2.3cm,bottom=2.2cm,left=1.8cm,right=1.8cm,headheight=1.3cm,headsep=0.4cm,footskip=1cm]{geometry}
\usepackage{amsmath,amssymb}
\usepackage{icomma}
\usepackage{graphicx}
\usepackage[table]{xcolor}
\usepackage{tikz}
\usepackage{fontawesome5}
\usepackage{ccicons}
\usepackage{fancyhdr}
\usepackage{lastpage}
\usepackage{tasks}
\usepackage{needspace}
\usepackage[most]{tcolorbox}
\usepackage{array}
\usepackage[hidelinks,pdftitle={Fitxa d'exercicis de la unitat 7: Geometria bàsica del pla. 1r d'ESO},pdfsubject={Matemàtiques 1r ESO},pdfkeywords={punt, recta, semirecta, segment, rectes perpendiculars, mediatriu, angles, bisectriu, sistema sexagesimal, complementaris, suplementaris, polígons, diagonals, triangles, quadrilàters, polígons regulars}]{hyperref}

\definecolor{navy}{HTML}{1B2748}
\definecolor{accent}{HTML}{E8603A}
\definecolor{teal}{HTML}{1F9E90}
\definecolor{soft}{HTML}{EEF1F8}
\definecolor{muted}{HTML}{5D6685}

\setlength{\parindent}{0pt}
\setlength{\parskip}{3pt}
\settasks{label=\alph*), label-format=\bfseries\color{navy}, label-width=1.3em, item-indent=1.8em, column-sep=0.6em, after-item-skip=5pt, before-skip=3pt, after-skip=4pt}

% ---------- capçalera i peu ----------
\newcommand{\llicencia}{\href{https://creativecommons.org/licenses/by-nc-sa/4.0/deed.ca}{\ccbyncsa}\ \ Llicència Creative Commons Reconeixement-NoComercial-CompartirIgual 4.0 Internacional (CC BY-NC-SA 4.0)}
\pagestyle{fancy}
\fancyhf{}
\fancyhead[L]{\raisebox{-0.35\height}{\includegraphics[height=0.95cm]{logo-gva.png}}}
\fancyhead[C]{\small\color{muted}Fitxa d'exercicis · Unitat 7 · 1r d'ESO}
\fancyhead[R]{\raisebox{-0.35\height}{\includegraphics[height=0.95cm]{logo-vermellar.png}}}
\fancyfoot[L]{\scriptsize\color{muted}\llicencia}
\fancyfoot[R]{\small\color{muted}\thepage}
\renewcommand{\headrulewidth}{0.4pt}
\renewcommand{\headrule}{\color{muted!40}\hrule height \headrulewidth}
\fancypagestyle{primera}{\fancyhf{}\renewcommand{\headrulewidth}{0pt}
  \fancyfoot[L]{\scriptsize\color{muted}\llicencia}
  \fancyfoot[R]{\small\color{muted}\thepage}}

% ---------- nivells i exercicis ----------
\newcommand{\estrelles}[1]{{\color{accent}\ifcase#1\or\faStar\faStar[regular]\faStar[regular]\or\faStar\faStar\faStar[regular]\or\faStar\faStar\faStar\fi}}
\newcounter{exer}
\newcommand{\ex}[1]{\par\needspace{4\baselineskip}\medskip\stepcounter{exer}\noindent{\color{navy}\bfseries\theexer.}\ \ {\small\estrelles{#1}}\ \ }

\newcommand{\seccio}[2]{\par\needspace{10\baselineskip}\bigskip
  \begin{tcolorbox}[colback=navy,colframe=navy,coltext=white,arc=2mm,boxrule=0pt,left=3mm,right=3mm,top=1.2mm,bottom=1.2mm]
  {\large\bfseries #1\quad #2}
  \end{tcolorbox}\vspace{-2pt}}

\newcolumntype{C}[1]{>{\centering\arraybackslash}p{#1}}

% quadre de recordatori (només a la fitxa d'exercicis)
\newcommand{\recorda}[1]{\par\smallskip\begin{tcolorbox}[colback=accent!7,colframe=accent!55,arc=2mm,boxrule=0.5pt,left=2.5mm,right=2.5mm,top=1mm,bottom=1mm]\small{\color{accent}\faLightbulb[regular]}\ \textbf{Recorda:}\ #1\end{tcolorbox}}
\newcommand{\res}[1]{#1}\newcommand{\rest}[1]{#1}

\def\fmneed{14}
% =====================================================================
% Capçalera FentMates de les fitxes: un color per a cada curs.
% Els colors són els de la pàgina inicial (LEVELS a assets/js/materials.js);
% el logo ve de _eines/logo/pdf/ (python3 _eines/logo/genera.py).
%
% Ús, al final del preàmbul de la unitat:
%   % =====================================================================
% Capçalera FentMates de les fitxes: un color per a cada curs.
% Els colors són els de la pàgina inicial (LEVELS a assets/js/materials.js);
% el logo ve de _eines/logo/pdf/ (python3 _eines/logo/genera.py).
%
% Ús, al final del preàmbul de la unitat:
%   % =====================================================================
% Capçalera FentMates de les fitxes: un color per a cada curs.
% Els colors són els de la pàgina inicial (LEVELS a assets/js/materials.js);
% el logo ve de _eines/logo/pdf/ (python3 _eines/logo/genera.py).
%
% Ús, al final del preàmbul de la unitat:
%   \input{../comu/fentmates.tex}
%   \fmcurs{1eso}
% i, al començament del document:
%   \capcalera{Matemàtiques}{1r ESO}{Fitxa d'exercicis · Unitat 2}{Divisibilitat}
% Opcional abans de \input: \def\fmneed{14} (espai mínim abans de cada secció).
% =====================================================================
\makeatletter
\def\fm@logodir{../../logo/pdf}
\providecommand{\fmneed}{10}
\def\fm@cursnom{}

% clau del curs, color i color del text damunt del color
\newcommand{\fm@curs}[3]{\@namedef{fm@color@#1}{#2}\@namedef{fm@text@#1}{#3}}
\fm@curs{1eso}{E6B84A}{navy}
\fm@curs{2eso}{3B82F6}{white}
\fm@curs{3eso}{FF7A50}{white}
\fm@curs{4eso-a}{A78BFA}{white}
\fm@curs{4eso-b}{6D28D9}{white}
\fm@curs{1batx-mi}{9B1C1C}{white}
\fm@curs{1batx-ccss}{EF6B6B}{white}
\fm@curs{2batx-mii}{1E6B47}{white}
\fm@curs{2batx-ccss}{2BB3A3}{white}

\newcommand{\fmcurs}[1]{%
  \@ifundefined{fm@color@#1}{\PackageError{fentmates}{Curs desconegut: #1}{}}{}%
  \edef\fm@tmp{\noexpand\definecolor{curs}{HTML}{\@nameuse{fm@color@#1}}%
               \noexpand\colorlet{curstext}{\@nameuse{fm@text@#1}}}\fm@tmp
  \colorlet{cursfosc}{curs!72!black}%
  \def\fm@logo{\fm@logodir/vertical-#1.pdf}%
  \def\fm@marca{\fm@logodir/marca-#1.pdf}}

% ---------- capçalera de la primera pàgina ----------
% #1 matèria (petit), #2 curs, #3 tipus de document i unitat, #4 títol
\newcommand{\capcalera}[4]{%
  \gdef\fm@cursnom{#2}%
  \par\noindent
  \begin{tcolorbox}[colback=curs,colframe=curs,boxrule=0pt,arc=2.2mm,width=3.5cm,height=2.15cm,
      valign=center,left=3mm,right=2mm,top=1mm,bottom=1mm,nobeforeafter,box align=center]
    \color{curstext}{\footnotesize #1}\par\vspace{1pt}{\fontsize{19}{21}\selectfont\bfseries #2\par}%
  \end{tcolorbox}\hspace{5mm}%
  \begin{minipage}[c]{\dimexpr\linewidth-3.5cm-5mm-2.9cm\relax}
    {\color{cursfosc}\small\bfseries\MakeUppercase{#3}}\par\vspace{2pt}
    {\color{navy}\fontsize{20}{23}\selectfont\bfseries #4\par}
  \end{minipage}\hfill
  \begin{minipage}[c]{2.6cm}\raggedleft\includegraphics[width=2.6cm]{\fm@logo}\end{minipage}\par
  \vspace{2pt}}

% ---------- barres de secció del color del curs ----------
\renewcommand{\seccio}[2]{\par\needspace{\fmneed\baselineskip}\bigskip
  \begin{tcolorbox}[colback=curs,colframe=curs,coltext=curstext,arc=2mm,boxrule=0pt,left=3mm,right=3mm,top=1.2mm,bottom=1.2mm]
  {\large\bfseries #1\quad #2}
  \end{tcolorbox}\vspace{-2pt}}

% ---------- estrelles de dificultat del color del curs (a Batxillerat no se'n posen) ----------
\@ifundefined{estrelles}{}{%
  \renewcommand{\estrelles}[1]{{\color{curs}\ifcase#1\or\faStar\faStar[regular]\faStar[regular]\or\faStar\faStar\faStar[regular]\or\faStar\faStar\faStar\fi}}}

% ---------- capçalera i peu de les pàgines ----------
% primera pàgina: logos institucionals al peu
\fancypagestyle{primera}{\fancyhf{}\renewcommand{\headrulewidth}{0pt}
  \fancyfoot[L]{\raisebox{-0.32\height}{\includegraphics[height=0.78cm]{logo-gva.png}}\hspace{5mm}%
    \raisebox{-0.32\height}{\includegraphics[height=0.78cm]{logo-vermellar.png}}}
  \fancyfoot[R]{\scriptsize\color{muted}\href{https://creativecommons.org/licenses/by-nc-sa/4.0/deed.ca}{\ccbyncsa}\ \ CC BY-NC-SA 4.0 · fentmates.org\qquad\small\thepage}}
% la resta: el símbol a l'esquerra i el curs a la dreta, amb el títol al centre com abans
\fancyhead[L]{\raisebox{-0.3\height}{\includegraphics[height=0.62cm]{\fm@marca}}}
\fancyhead[R]{\small\bfseries\color{cursfosc}\fm@cursnom}
\renewcommand{\headrule}{\color{curs!45}\hrule height \headrulewidth}
\makeatother

%   \fmcurs{1eso}
% i, al començament del document:
%   \capcalera{Matemàtiques}{1r ESO}{Fitxa d'exercicis · Unitat 2}{Divisibilitat}
% Opcional abans de \input: \def\fmneed{14} (espai mínim abans de cada secció).
% =====================================================================
\makeatletter
\def\fm@logodir{../../logo/pdf}
\providecommand{\fmneed}{10}
\def\fm@cursnom{}

% clau del curs, color i color del text damunt del color
\newcommand{\fm@curs}[3]{\@namedef{fm@color@#1}{#2}\@namedef{fm@text@#1}{#3}}
\fm@curs{1eso}{E6B84A}{navy}
\fm@curs{2eso}{3B82F6}{white}
\fm@curs{3eso}{FF7A50}{white}
\fm@curs{4eso-a}{A78BFA}{white}
\fm@curs{4eso-b}{6D28D9}{white}
\fm@curs{1batx-mi}{9B1C1C}{white}
\fm@curs{1batx-ccss}{EF6B6B}{white}
\fm@curs{2batx-mii}{1E6B47}{white}
\fm@curs{2batx-ccss}{2BB3A3}{white}

\newcommand{\fmcurs}[1]{%
  \@ifundefined{fm@color@#1}{\PackageError{fentmates}{Curs desconegut: #1}{}}{}%
  \edef\fm@tmp{\noexpand\definecolor{curs}{HTML}{\@nameuse{fm@color@#1}}%
               \noexpand\colorlet{curstext}{\@nameuse{fm@text@#1}}}\fm@tmp
  \colorlet{cursfosc}{curs!72!black}%
  \def\fm@logo{\fm@logodir/vertical-#1.pdf}%
  \def\fm@marca{\fm@logodir/marca-#1.pdf}}

% ---------- capçalera de la primera pàgina ----------
% #1 matèria (petit), #2 curs, #3 tipus de document i unitat, #4 títol
\newcommand{\capcalera}[4]{%
  \gdef\fm@cursnom{#2}%
  \par\noindent
  \begin{tcolorbox}[colback=curs,colframe=curs,boxrule=0pt,arc=2.2mm,width=3.5cm,height=2.15cm,
      valign=center,left=3mm,right=2mm,top=1mm,bottom=1mm,nobeforeafter,box align=center]
    \color{curstext}{\footnotesize #1}\par\vspace{1pt}{\fontsize{19}{21}\selectfont\bfseries #2\par}%
  \end{tcolorbox}\hspace{5mm}%
  \begin{minipage}[c]{\dimexpr\linewidth-3.5cm-5mm-2.9cm\relax}
    {\color{cursfosc}\small\bfseries\MakeUppercase{#3}}\par\vspace{2pt}
    {\color{navy}\fontsize{20}{23}\selectfont\bfseries #4\par}
  \end{minipage}\hfill
  \begin{minipage}[c]{2.6cm}\raggedleft\includegraphics[width=2.6cm]{\fm@logo}\end{minipage}\par
  \vspace{2pt}}

% ---------- barres de secció del color del curs ----------
\renewcommand{\seccio}[2]{\par\needspace{\fmneed\baselineskip}\bigskip
  \begin{tcolorbox}[colback=curs,colframe=curs,coltext=curstext,arc=2mm,boxrule=0pt,left=3mm,right=3mm,top=1.2mm,bottom=1.2mm]
  {\large\bfseries #1\quad #2}
  \end{tcolorbox}\vspace{-2pt}}

% ---------- estrelles de dificultat del color del curs (a Batxillerat no se'n posen) ----------
\@ifundefined{estrelles}{}{%
  \renewcommand{\estrelles}[1]{{\color{curs}\ifcase#1\or\faStar\faStar[regular]\faStar[regular]\or\faStar\faStar\faStar[regular]\or\faStar\faStar\faStar\fi}}}

% ---------- capçalera i peu de les pàgines ----------
% primera pàgina: logos institucionals al peu
\fancypagestyle{primera}{\fancyhf{}\renewcommand{\headrulewidth}{0pt}
  \fancyfoot[L]{\raisebox{-0.32\height}{\includegraphics[height=0.78cm]{logo-gva.png}}\hspace{5mm}%
    \raisebox{-0.32\height}{\includegraphics[height=0.78cm]{logo-vermellar.png}}}
  \fancyfoot[R]{\scriptsize\color{muted}\href{https://creativecommons.org/licenses/by-nc-sa/4.0/deed.ca}{\ccbyncsa}\ \ CC BY-NC-SA 4.0 · fentmates.org\qquad\small\thepage}}
% la resta: el símbol a l'esquerra i el curs a la dreta, amb el títol al centre com abans
\fancyhead[L]{\raisebox{-0.3\height}{\includegraphics[height=0.62cm]{\fm@marca}}}
\fancyhead[R]{\small\bfseries\color{cursfosc}\fm@cursnom}
\renewcommand{\headrule}{\color{curs!45}\hrule height \headrulewidth}
\makeatother

%   \fmcurs{1eso}
% i, al començament del document:
%   \capcalera{Matemàtiques}{1r ESO}{Fitxa d'exercicis · Unitat 2}{Divisibilitat}
% Opcional abans de \input: \def\fmneed{14} (espai mínim abans de cada secció).
% =====================================================================
\makeatletter
\def\fm@logodir{../../logo/pdf}
\providecommand{\fmneed}{10}
\def\fm@cursnom{}

% clau del curs, color i color del text damunt del color
\newcommand{\fm@curs}[3]{\@namedef{fm@color@#1}{#2}\@namedef{fm@text@#1}{#3}}
\fm@curs{1eso}{E6B84A}{navy}
\fm@curs{2eso}{3B82F6}{white}
\fm@curs{3eso}{FF7A50}{white}
\fm@curs{4eso-a}{A78BFA}{white}
\fm@curs{4eso-b}{6D28D9}{white}
\fm@curs{1batx-mi}{9B1C1C}{white}
\fm@curs{1batx-ccss}{EF6B6B}{white}
\fm@curs{2batx-mii}{1E6B47}{white}
\fm@curs{2batx-ccss}{2BB3A3}{white}

\newcommand{\fmcurs}[1]{%
  \@ifundefined{fm@color@#1}{\PackageError{fentmates}{Curs desconegut: #1}{}}{}%
  \edef\fm@tmp{\noexpand\definecolor{curs}{HTML}{\@nameuse{fm@color@#1}}%
               \noexpand\colorlet{curstext}{\@nameuse{fm@text@#1}}}\fm@tmp
  \colorlet{cursfosc}{curs!72!black}%
  \def\fm@logo{\fm@logodir/vertical-#1.pdf}%
  \def\fm@marca{\fm@logodir/marca-#1.pdf}}

% ---------- capçalera de la primera pàgina ----------
% #1 matèria (petit), #2 curs, #3 tipus de document i unitat, #4 títol
\newcommand{\capcalera}[4]{%
  \gdef\fm@cursnom{#2}%
  \par\noindent
  \begin{tcolorbox}[colback=curs,colframe=curs,boxrule=0pt,arc=2.2mm,width=3.5cm,height=2.15cm,
      valign=center,left=3mm,right=2mm,top=1mm,bottom=1mm,nobeforeafter,box align=center]
    \color{curstext}{\footnotesize #1}\par\vspace{1pt}{\fontsize{19}{21}\selectfont\bfseries #2\par}%
  \end{tcolorbox}\hspace{5mm}%
  \begin{minipage}[c]{\dimexpr\linewidth-3.5cm-5mm-2.9cm\relax}
    {\color{cursfosc}\small\bfseries\MakeUppercase{#3}}\par\vspace{2pt}
    {\color{navy}\fontsize{20}{23}\selectfont\bfseries #4\par}
  \end{minipage}\hfill
  \begin{minipage}[c]{2.6cm}\raggedleft\includegraphics[width=2.6cm]{\fm@logo}\end{minipage}\par
  \vspace{2pt}}

% ---------- barres de secció del color del curs ----------
\renewcommand{\seccio}[2]{\par\needspace{\fmneed\baselineskip}\bigskip
  \begin{tcolorbox}[colback=curs,colframe=curs,coltext=curstext,arc=2mm,boxrule=0pt,left=3mm,right=3mm,top=1.2mm,bottom=1.2mm]
  {\large\bfseries #1\quad #2}
  \end{tcolorbox}\vspace{-2pt}}

% ---------- estrelles de dificultat del color del curs (a Batxillerat no se'n posen) ----------
\@ifundefined{estrelles}{}{%
  \renewcommand{\estrelles}[1]{{\color{curs}\ifcase#1\or\faStar\faStar[regular]\faStar[regular]\or\faStar\faStar\faStar[regular]\or\faStar\faStar\faStar\fi}}}

% ---------- capçalera i peu de les pàgines ----------
% primera pàgina: logos institucionals al peu
\fancypagestyle{primera}{\fancyhf{}\renewcommand{\headrulewidth}{0pt}
  \fancyfoot[L]{\raisebox{-0.32\height}{\includegraphics[height=0.78cm]{logo-gva.png}}\hspace{5mm}%
    \raisebox{-0.32\height}{\includegraphics[height=0.78cm]{logo-vermellar.png}}}
  \fancyfoot[R]{\scriptsize\color{muted}\href{https://creativecommons.org/licenses/by-nc-sa/4.0/deed.ca}{\ccbyncsa}\ \ CC BY-NC-SA 4.0 · fentmates.org\qquad\small\thepage}}
% la resta: el símbol a l'esquerra i el curs a la dreta, amb el títol al centre com abans
\fancyhead[L]{\raisebox{-0.3\height}{\includegraphics[height=0.62cm]{\fm@marca}}}
\fancyhead[R]{\small\bfseries\color{cursfosc}\fm@cursnom}
\renewcommand{\headrule}{\color{curs!45}\hrule height \headrulewidth}
\makeatother

\fmcurs{1eso}

\begin{document}
\thispagestyle{primera}

\capcalera{Matemàtiques}{1r ESO}{Fitxa d'exercicis · Unitat 7}{Geometria bàsica del pla}
\vspace{10pt}
Nom i cognoms: \rule{8.4cm}{0.4pt}\hfill Grup: \rule{1.6cm}{0.4pt}\hfill Data: \rule{2.2cm}{0.4pt}
\vspace{8pt}

\begin{tcolorbox}[colback=soft,colframe=soft,arc=2mm,left=3mm,right=3mm,top=1.5mm,bottom=1.5mm]
\small Exercicis graduats per nivells de dificultat:\quad \estrelles{1}\ bàsic \qquad \estrelles{2}\ mitjà \qquad \estrelles{3}\ avançat.\\[2pt] Resol-los al quadern mostrant tots els passos. Tin a mà el regle, el compàs i el transportador d'angles, i fes un dibuix sempre que t'ajude.
\end{tcolorbox}

% =====================================================
\seccio{1}{Punts, rectes, semirectes i segments}

\ex{1} Observa la figura i respon:\par\smallskip\begin{center}\begin{tikzpicture}[x=0.6cm,y=0.6cm,baseline=(current bounding box.center),line join=round,line cap=round]
\draw[thick,latex-latex] (-0.500,0.000) -- (7.500,0.000);
\node[font=\small,below] at (7.300,0.000) {$r$};
\fill (1.000,0.000) circle (1.5pt);
\node[font=\small,above] at (1.000,0.000) {$A$};
\fill (3.500,0.000) circle (1.5pt);
\node[font=\small,above] at (3.500,0.000) {$B$};
\fill (6.000,0.000) circle (1.5pt);
\node[font=\small,above] at (6.000,0.000) {$C$};
\fill (4.600,1.500) circle (1.5pt);
\node[font=\small,above] at (4.600,1.500) {$D$};
\end{tikzpicture}\end{center}\begin{tasks}(1)
  \task Quins punts són de la recta $r$? Quin punt no hi és?
  \task Escriu tots els segments que tenen els extrems en dos dels punts $A$, $B$ i $C$.
  \task Quantes semirectes de la recta $r$ tenen l'origen en $B$?
\end{tasks}\recorda{La \textbf{recta} no té principi ni final; la \textbf{semirecta} té origen però no té final; el \textbf{segment} té dos extrems i es pot mesurar. Per dos punts passa una única recta.}

\ex{1} Digues si cada figura és una recta, una semirecta o un segment:\begin{tasks}(3)
  \task \begin{tikzpicture}[x=0.6cm,y=0.6cm,baseline=(current bounding box.center),line join=round,line cap=round]
\draw[thick,latex-latex] (0.000,0.000) -- (3.000,0.600);
\fill (1.000,0.200) circle (1.5pt);
\fill (2.000,0.400) circle (1.5pt);
\end{tikzpicture}
  \task \begin{tikzpicture}[x=0.6cm,y=0.6cm,baseline=(current bounding box.center),line join=round,line cap=round]
\draw[thick,-] (0.000,0.000) -- (3.000,0.600);
\fill (0.000,0.000) circle (1.5pt);
\fill (3.000,0.600) circle (1.5pt);
\end{tikzpicture}
  \task \begin{tikzpicture}[x=0.6cm,y=0.6cm,baseline=(current bounding box.center),line join=round,line cap=round]
\draw[thick,-latex] (0.000,0.000) -- (3.000,0.600);
\fill (0.000,0.000) circle (1.5pt);
\end{tikzpicture}
\end{tasks}

\ex{1} Respon:\begin{tasks}(1)
  \task El segment $\overline{AB}$ mesura 7\,cm i $M$ és el seu punt mitjà. Quant mesura $\overline{AM}$?
  \task $M$ és el punt mitjà de $\overline{PQ}$ i $\overline{PM}$ fa 4,2\,cm. Quant mesura $\overline{PQ}$?
\end{tasks}

\ex{2} Sobre una recta marquem quatre punts. Quants segments determinen? I si en marquem sis?

\ex{2} Els punts $A$, $B$ i $C$ estan sobre una recta, en aquest ordre. $\overline{AB}=5$\,cm i $\overline{BC}=3$\,cm. Quant mesura $\overline{AC}$? I si $C$ estiguera entre $A$ i $B$ (amb $\overline{BC}=3$\,cm)?

\ex{3} Dibuixem punts de manera que no n'hi ha tres d'alineats. Quantes rectes podem traçar que passen per dos d'ells si hi ha 4 punts? I 5? I 7?


% =====================================================
\seccio{2}{Posicions relatives de rectes. Mediatriu}

\ex{1} Les rectes $r$ i $s$ són paral·leles. Digues com són:\par\smallskip\begin{center}\begin{tikzpicture}[x=0.6cm,y=0.6cm,baseline=(current bounding box.center),line join=round,line cap=round]
\draw[thick,latex-latex] (-0.400,0.000) -- (6.400,0.000);
\draw[thick,latex-latex] (-0.400,1.600) -- (6.400,1.600);
\draw[thick,latex-latex] (1.500,-0.900) -- (1.500,2.600);
\draw[thick,latex-latex] (2.800,-0.900) -- (6.200,2.600);
\draw[thin] (1.867,0.000) -- (1.867,0.367) -- (1.500,0.367);
\node[font=\small,right] at (6.400,0.000) {$r$};
\node[font=\small,right] at (6.400,1.600) {$s$};
\node[font=\small,above] at (1.500,2.600) {$t$};
\node[font=\small,above] at (6.200,2.600) {$u$};
\end{tikzpicture}\end{center}\begin{tasks}(4)
  \task $r$ i $s$
  \task $r$ i $t$
  \task $s$ i $t$
  \task $r$ i $u$
\end{tasks}\recorda{Dues rectes són \textbf{secants} si es tallen en un punt, \textbf{paral·leles} si no es tallen mai i \textbf{perpendiculars} si es tallen formant quatre angles rectes ($r\perp s$).}

\ex{1} Digues si són paral·leles, secants o perpendiculars:\begin{tasks}(2)
  \task Dos costats oposats d'un full de paper.
  \task Dos costats d'un full que formen una cantonada.
  \task Les agulles d'un rellotge a les 9 en punt.
  \task Les agulles d'un rellotge a les 10 en punt.
\end{tasks}

\ex{1} Dibuixa un segment $\overline{AB}$ de 6\,cm i traça'n la mediatriu amb regle i compàs. A quina distància d'$A$ talla el segment?

\ex{2} Digues si és cert o fals:\begin{tasks}(1)
  \task Dues rectes perpendiculars són sempre secants.
  \task Dues rectes secants són sempre perpendiculars.
  \task Si $r\parallel s$ i $s\parallel t$, aleshores $r\parallel t$.
  \task Si $r\perp s$ i $s\perp t$, aleshores $r\perp t$.
\end{tasks}

\ex{2} El punt $P$ està sobre la mediatriu del segment $\overline{AB}$, que mesura 8\,cm, i la distància de $P$ a $A$ és 5\,cm. Quant mesura $\overline{PB}$? A quina distància de $A$ està el punt mitjà $M$?

\ex{3} Dues cases, $A$ i $B$, estan a la mateixa banda d'un camí recte. On s'ha de posar una font, vora el camí, perquè estiga a la mateixa distància de les dues cases? Explica com la trobaries.


% =====================================================
\seccio{3}{Angles: elements, mesura i classificació}

\ex{1} Classifica aquests angles (nul, agut, recte, obtús, pla, còncau o complet):\begin{tasks}(4)
  \task $25^\circ$
  \task $90^\circ$
  \task $100^\circ$
  \task $180^\circ$
  \task $270^\circ$
  \task $360^\circ$
  \task $179^\circ$
  \task $89^\circ$
\end{tasks}\recorda{Agut: menys de $90^\circ$. Recte: $90^\circ$. Obtús: entre $90^\circ$ i $180^\circ$. Pla: $180^\circ$. Còncau: més de $180^\circ$. Complet: $360^\circ$.}

\ex{1} Classifica cada angle i mesura'l amb el transportador:\begin{tasks}(2)
  \task \begin{tikzpicture}[x=0.8cm,y=0.8cm,baseline=(current bounding box.center),line join=round,line cap=round]
\draw[fill=accent!30,draw=accent] (0.000,0.000) -- (0.591,0.104) arc[start angle=10,end angle=50,radius=0.6] -- cycle;
\draw[thick,-latex] (0.000,0.000) -- (2.954,0.521);
\draw[thick,-latex] (0.000,0.000) -- (1.928,2.298);
\fill (0.000,0.000) circle (1.5pt);
\end{tikzpicture}
  \task \begin{tikzpicture}[x=0.8cm,y=0.8cm,baseline=(current bounding box.center),line join=round,line cap=round]
\draw[fill=accent!30,draw=accent] (0.000,0.000) -- (0.400,0.000) -- (0.400,0.400) -- (0.000,0.400) -- cycle;
\draw[thick,-latex] (0.000,0.000) -- (3.000,0.000);
\draw[thick,-latex] (0.000,0.000) -- (0.000,3.000);
\fill (0.000,0.000) circle (1.5pt);
\end{tikzpicture}
  \task \begin{tikzpicture}[x=0.8cm,y=0.8cm,baseline=(current bounding box.center),line join=round,line cap=round]
\draw[fill=accent!30,draw=accent] (0.000,0.000) -- (0.600,0.000) arc[start angle=0,end angle=140,radius=0.6] -- cycle;
\draw[thick,-latex] (0.000,0.000) -- (3.000,0.000);
\draw[thick,-latex] (0.000,0.000) -- (-2.298,1.928);
\fill (0.000,0.000) circle (1.5pt);
\end{tikzpicture}
  \task \begin{tikzpicture}[x=0.8cm,y=0.8cm,baseline=(current bounding box.center),line join=round,line cap=round]
\draw[fill=accent!30,draw=accent] (0.000,0.000) -- (0.600,0.000) arc[start angle=0,end angle=230,radius=0.6] -- cycle;
\draw[thick,-latex] (0.000,0.000) -- (3.000,0.000);
\draw[thick,-latex] (0.000,0.000) -- (-1.928,-2.298);
\fill (0.000,0.000) circle (1.5pt);
\end{tikzpicture}
\end{tasks}

\ex{1} Quin angle (convex) formen les agulles d'un rellotge a les hores en punt següents?\begin{tasks}(6)
  \task 3 h
  \task 4 h
  \task 5 h
  \task 6 h
  \task 9 h
  \task 11 h
\end{tasks}

\ex{2} Quants graus gira l'agulla dels minuts en aquests temps?\begin{tasks}(4)
  \task 10 min
  \task 25 min
  \task 1 hora
  \task tres quarts d'hora
\end{tasks}

\ex{2} Respon:\begin{tasks}(1)
  \task Quins angles forma la bisectriu d'un angle de $136^\circ$ amb cadascun dels costats?
  \task La bisectriu d'un angle forma amb un dels costats un angle de $43^\circ$. Quant mesura l'angle?
\end{tasks}

\ex{2} Dues semirectes amb el mateix origen formen un angle convex de $125^\circ$. Quant mesura l'angle còncau que formen?

\ex{3} Respon:\begin{tasks}(1)
  \task Dividim un angle complet en 5 angles iguals. Quant mesura cadascun? I si el dividim en 8?
  \task Quantes vegades cap un angle de $15^\circ$ en un angle pla?
\end{tasks}


% =====================================================
\seccio{4}{Mesura d'angles: el sistema sexagesimal}

\ex{1} Expressa en segons:\begin{tasks}(3)
  \task $2^\circ$
  \task $15'$
  \task $1^\circ\,20'$
  \task $3^\circ\,5'\,40''$
  \task $10^\circ\,10'\,10''$
  \task $45'\,30''$
\end{tasks}\recorda{$1^\circ=60'$, \ $1'=60''$, \ $1^\circ=3600''$. Per a passar a una unitat més menuda multipliquem; per a passar a una unitat més gran dividim entre 60.}

\ex{1} Expressa en minuts:\begin{tasks}(4)
  \task $5^\circ$
  \task $2^\circ\,30'$
  \task $12^\circ\,45'$
  \task $180^\circ$
\end{tasks}

\ex{1} Passa a forma complexa (graus, minuts i segons):\begin{tasks}(3)
  \task $150'$
  \task $4000''$
  \task $7385''$
  \task $10\,000''$
  \task $1000'$
  \task $86\,400''$
\end{tasks}

\ex{2} Expressa en graus i minuts:\begin{tasks}(4)
  \task $12,5^\circ$
  \task $30,25^\circ$
  \task $8,75^\circ$
  \task $0,1^\circ$
\end{tasks}

\ex{2} Calcula:\begin{tasks}(3)
  \task $23^\circ\,45'+18^\circ\,30'$
  \task $45^\circ\,36'\,50''+27^\circ\,48'\,25''$
  \task $120^\circ\,15'\,40''+59^\circ\,44'\,20''$
\end{tasks}

\ex{2} Calcula:\begin{tasks}(3)
  \task $75^\circ\,40'-32^\circ\,15'$
  \task $60^\circ\,20'\,15''-25^\circ\,45'\,30''$
  \task $180^\circ-97^\circ\,34'\,18''$
\end{tasks}

\ex{2} Calcula:\begin{tasks}(3)
  \task $(12^\circ\,20')\cdot3$
  \task $(25^\circ\,15'\,40'')\cdot4$
  \task $(8^\circ\,45'\,50'')\cdot6$
\end{tasks}

\ex{3} Calcula i comprova el resultat multiplicant:\begin{tasks}(3)
  \task $(75^\circ\,36'):4$
  \task $(100^\circ\,20'\,30''):5$
  \task $(47^\circ\,13'):3$
\end{tasks}

\ex{3} La Terra fa una volta completa sobre ella mateixa en 24 hores. Quants graus gira en 1 hora? I en 1 minut? I en 1 segon?


% =====================================================
\seccio{5}{Relacions entre angles}

\ex{1} Calcula el complementari de cada angle:\begin{tasks}(4)
  \task $30^\circ$
  \task $47^\circ$
  \task $89^\circ$
  \task $15^\circ\,30'$
\end{tasks}\recorda{Dos angles són \textbf{complementaris} si sumen $90^\circ$ i \textbf{suplementaris} si sumen $180^\circ$. Els angles \textbf{oposats pel vèrtex} són iguals i els \textbf{adjacents} sumen $180^\circ$.}

\ex{1} Calcula el suplementari de cada angle:\begin{tasks}(4)
  \task $45^\circ$
  \task $120^\circ$
  \task $179^\circ$
  \task $92^\circ\,15'\,40''$
\end{tasks}

\ex{1} Dues rectes es tallen i un dels angles fa $35^\circ$. Calcula els angles $a$, $b$ i $c$ i explica com els has trobat.\par\smallskip\begin{center}\begin{tikzpicture}[x=0.6cm,y=0.6cm,baseline=(current bounding box.center),line join=round,line cap=round]
\draw[fill=accent!30,draw=accent] (0.000,0.000) -- (0.433,0.250) arc[start angle=30,end angle=65,radius=0.5] -- cycle;
\draw[thick,latex-latex] (-1.732,-1.000) -- (1.732,1.000);
\draw[thick,latex-latex] (-0.845,-1.813) -- (0.845,1.813);
\fill (0.000,0.000) circle (1.5pt);
\node[font=\small,inner sep=1pt] at (0.648,0.695) {$35^\circ$};
\node[font=\small,inner sep=1pt] at (-0.629,0.408) {$a$};
\node[font=\small,inner sep=1pt] at (-0.511,-0.549) {$b$};
\node[font=\small,inner sep=1pt] at (0.629,-0.408) {$c$};
\end{tikzpicture}\end{center}

\ex{2} Les rectes $r$ i $s$ són paral·leles i l'angle 1 fa $58^\circ$. Calcula la resta d'angles numerats.\par\smallskip\begin{center}\begin{tikzpicture}[x=0.85cm,y=0.85cm,baseline=(current bounding box.center),line join=round,line cap=round]
\draw[thick,latex-latex] (-2.200,0.000) -- (3.325,0.000);
\draw[thick,latex-latex] (-2.200,1.800) -- (3.325,1.800);
\draw[thick,latex-latex] (-0.689,-1.102) -- (1.814,2.902);
\node[font=\small,above] at (-2.200,1.800) {$r$};
\node[font=\small,above] at (-2.200,0.000) {$s$};
\node[font=\small,inner sep=1pt] at (0.542,0.301) {\scriptsize 1};
\node[font=\small,inner sep=1pt] at (-0.301,0.542) {\scriptsize 2};
\node[font=\small,inner sep=1pt] at (-0.542,-0.301) {\scriptsize 3};
\node[font=\small,inner sep=1pt] at (0.301,-0.542) {\scriptsize 4};
\node[font=\small,inner sep=1pt] at (1.667,2.101) {\scriptsize 5};
\node[font=\small,inner sep=1pt] at (0.824,2.342) {\scriptsize 6};
\node[font=\small,inner sep=1pt] at (0.583,1.499) {\scriptsize 7};
\node[font=\small,inner sep=1pt] at (1.425,1.258) {\scriptsize 8};
\end{tikzpicture}\end{center}

\ex{2} Digues si és cert o fals:\begin{tasks}(1)
  \task Dos angles aguts poden ser suplementaris.
  \task El suplementari d'un angle agut és obtús.
  \task Dos angles adjacents sempre són suplementaris.
  \task Un angle obtús té complementari.
\end{tasks}

\ex{2} Calcula el complementari i el suplementari de l'angle $37^\circ\,25'\,48''$.

\ex{3} Dos angles són complementaris i un d'ells mesura $20^\circ$ més que l'altre. Quant mesura cadascun?

\ex{3} El suplementari d'un angle és el quàdruple de l'angle. Quant mesura l'angle?

\ex{3} Dues rectes paral·leles tallades per una secant formen dos angles alterns interns que mesuren $3x+10^\circ$ i $70^\circ$. Quant val $x$?


% =====================================================
\seccio{6}{Polígons: elements i classificació}

\ex{1} Escriu el nom del polígon que té:\begin{tasks}(3)
  \task 5 costats
  \task 8 costats
  \task 6 costats
  \task 10 costats
  \task 7 costats
  \task 12 costats
\end{tasks}

\ex{1} Per a cada polígon, digues quants costats té, si és convex o còncau i si és regular o irregular:\begin{tasks}(4)
  \task \begin{tikzpicture}[x=0.6cm,y=0.6cm,baseline=(current bounding box.center),line join=round,line cap=round]
\draw[fill=teal!12,draw=teal,thick] (-0.647,-0.890) -- (0.647,-0.890) -- (1.046,0.340) -- (0.000,1.100) -- (-1.046,0.340) -- cycle;
\end{tikzpicture}
  \task \begin{tikzpicture}[x=0.6cm,y=0.6cm,baseline=(current bounding box.center),line join=round,line cap=round]
\draw[fill=teal!12,draw=teal,thick] (0.000,0.000) -- (2.400,0.000) -- (2.400,2.000) -- (1.200,0.900) -- (0.000,2.000) -- cycle;
\end{tikzpicture}
  \task \begin{tikzpicture}[x=0.6cm,y=0.6cm,baseline=(current bounding box.center),line join=round,line cap=round]
\draw[fill=teal!12,draw=teal,thick] (0.000,0.000) -- (2.800,0.000) -- (2.800,1.600) -- (0.000,1.600) -- cycle;
\end{tikzpicture}
  \task \begin{tikzpicture}[x=0.6cm,y=0.6cm,baseline=(current bounding box.center),line join=round,line cap=round]
\draw[fill=teal!12,draw=teal,thick] (0.000,0.000) -- (2.400,0.300) -- (2.000,1.900) -- (0.300,1.500) -- cycle;
\end{tikzpicture}
\end{tasks}\recorda{Un polígon és \textbf{còncau} si té algun angle interior de més de $180^\circ$, i \textbf{regular} si té tots els costats \emph{i} tots els angles iguals.}

\ex{1} Quantes diagonals ixen de cada vèrtex i quantes en té en total?\begin{tasks}(4)
  \task quadrilàter
  \task pentàgon
  \task octàgon
  \task decàgon
\end{tasks}

\ex{2} Digues si és cert o fals:\begin{tasks}(1)
  \task Un polígon regular té tots els costats iguals.
  \task Un rombe (4 costats iguals) és un polígon regular.
  \task Un polígon còncau té algun angle interior de més de $180^\circ$.
  \task Un triangle té 3 diagonals.
\end{tasks}

\ex{2} Quin polígon té 14 diagonals? I quin en té 27?

\ex{3} De cada vèrtex d'un polígon ixen 9 diagonals. Quants costats té? Quantes diagonals té en total?


% =====================================================
\seccio{7}{Triangles}

\ex{1} Classifica segons els costats els triangles que tenen aquestes mesures (en cm):\begin{tasks}(4)
  \task 5, 5, 5
  \task 4, 7, 7
  \task 3, 5, 6
  \task 8, 8, 3
\end{tasks}\recorda{Segons els costats: \textbf{equilàter} (3 iguals), \textbf{isòsceles} (2 iguals), \textbf{escalé} (cap igual). Segons els angles: \textbf{acutangle}, \textbf{rectangle} o \textbf{obtusangle}. Els tres angles sumen $180^\circ$.}

\ex{1} Dos angles d'un triangle mesuren el que s'indica. Calcula el tercer i classifica el triangle segons els angles:\begin{tasks}(4)
  \task $40^\circ$ i $60^\circ$
  \task $90^\circ$ i $35^\circ$
  \task $120^\circ$ i $25^\circ$
  \task $45^\circ$ i $45^\circ$
\end{tasks}

\ex{1} Es pot construir un triangle amb segments d'aquestes longituds (en cm)?\begin{tasks}(4)
  \task 3, 4, 5
  \task 2, 3, 6
  \task 5, 5, 9
  \task 4, 4, 8
\end{tasks}

\ex{2} Calcula:\begin{tasks}(1)
  \task Els angles de la base d'un triangle isòsceles en què l'angle desigual fa $50^\circ$.
  \task L'angle desigual d'un triangle isòsceles en què cada angle de la base fa $72^\circ$.
  \task Els angles aguts d'un triangle rectangle isòsceles.
\end{tasks}

\ex{2} Pot existir un triangle que siga…\begin{tasks}(2)
  \task rectangle i isòsceles?
  \task obtusangle i equilàter?
  \task acutangle i escalé?
  \task rectangle i equilàter?
\end{tasks}

\ex{2} Un angle agut d'un triangle rectangle mesura $27^\circ\,35'$. Quant mesura l'altre angle agut?

\ex{2} Quin punt notable d'un triangle és…\begin{tasks}(1)
  \task el centre de la circumferència que passa pels tres vèrtexs?
  \task el punt on es tallen les tres altures?
  \task el centre de gravetat del triangle?
  \task el punt que està a la mateixa distància dels tres costats?
\end{tasks}

\ex{3} Els angles d'un triangle mesuren $x$, $x+15^\circ$ i $x+30^\circ$. Calcula'ls i classifica el triangle.

\ex{3} Dos costats d'un triangle mesuren 5\,cm i 12\,cm. Quines longituds (en centímetres sencers) pot tindre el tercer costat?

\ex{3} Prolonguem un costat d'un triangle i, en un vèrtex, l'angle interior i l'angle adjacent que es forma fora del triangle fan $130^\circ$ (l'adjacent). Un altre angle del triangle fa $60^\circ$. Calcula els tres angles del triangle.


% =====================================================
\seccio{8}{Quadrilàters}

\ex{1} Escriu el nom de cada quadrilàter (les marques indiquen costats iguals):\begin{tasks}(3)
  \task \begin{tikzpicture}[x=0.5cm,y=0.5cm,baseline=(current bounding box.center),line join=round,line cap=round]
\draw[fill=teal!12,draw=teal,thick] (0.000,0.000) -- (2.000,0.000) -- (2.000,2.000) -- (0.000,2.000) -- cycle;
\draw[thin] (0.000,0.440) -- (0.440,0.440) -- (0.440,0.000);
\draw[thin] (1.560,0.000) -- (1.560,0.440) -- (2.000,0.440);
\draw[thin] (2.000,1.560) -- (1.560,1.560) -- (1.560,2.000);
\draw[thin] (0.440,2.000) -- (0.440,1.560) -- (0.000,1.560);
\draw[thick] (1.000,0.240) -- (1.000,-0.240);
\draw[thick] (1.760,1.000) -- (2.240,1.000);
\draw[thick] (1.000,1.760) -- (1.000,2.240);
\draw[thick] (0.240,1.000) -- (-0.240,1.000);
\end{tikzpicture}
  \task \begin{tikzpicture}[x=0.5cm,y=0.5cm,baseline=(current bounding box.center),line join=round,line cap=round]
\draw[fill=teal!12,draw=teal,thick] (0.000,0.000) -- (1.600,-1.100) -- (3.200,0.000) -- (1.600,1.100) -- cycle;
\draw[thick] (0.936,-0.352) -- (0.664,-0.748);
\draw[thick] (2.264,-0.352) -- (2.536,-0.748);
\draw[thick] (2.264,0.352) -- (2.536,0.748);
\draw[thick] (0.936,0.352) -- (0.664,0.748);
\end{tikzpicture}
  \task \begin{tikzpicture}[x=0.5cm,y=0.5cm,baseline=(current bounding box.center),line join=round,line cap=round]
\draw[fill=teal!12,draw=teal,thick] (0.000,0.000) -- (3.400,0.000) -- (2.500,1.700) -- (0.900,1.700) -- cycle;
\draw[thick] (2.738,0.738) -- (3.162,0.962);
\draw[thick] (0.662,0.738) -- (0.238,0.962);
\end{tikzpicture}
  \task \begin{tikzpicture}[x=0.5cm,y=0.5cm,baseline=(current bounding box.center),line join=round,line cap=round]
\draw[fill=teal!12,draw=teal,thick] (0.000,0.000) -- (3.000,0.000) -- (3.000,1.700) -- (0.000,1.700) -- cycle;
\draw[thin] (0.000,0.440) -- (0.440,0.440) -- (0.440,0.000);
\draw[thin] (2.560,0.000) -- (2.560,0.440) -- (3.000,0.440);
\draw[thin] (3.000,1.260) -- (2.560,1.260) -- (2.560,1.700);
\draw[thin] (0.440,1.700) -- (0.440,1.260) -- (0.000,1.260);
\draw[thick] (1.500,0.240) -- (1.500,-0.240);
\draw[thick] (2.760,0.760) -- (3.240,0.760);
\draw[thick] (2.760,0.940) -- (3.240,0.940);
\draw[thick] (1.500,1.460) -- (1.500,1.940);
\draw[thick] (0.240,0.940) -- (-0.240,0.940);
\draw[thick] (0.240,0.760) -- (-0.240,0.760);
\end{tikzpicture}
  \task \begin{tikzpicture}[x=0.5cm,y=0.5cm,baseline=(current bounding box.center),line join=round,line cap=round]
\draw[fill=teal!12,draw=teal,thick] (0.000,0.000) -- (2.600,0.000) -- (3.500,1.600) -- (0.900,1.600) -- cycle;
\draw[thick] (1.300,0.240) -- (1.300,-0.240);
\draw[thick] (2.797,0.839) -- (3.215,0.604);
\draw[thick] (2.885,0.996) -- (3.303,0.761);
\draw[thick] (2.200,1.360) -- (2.200,1.840);
\draw[thick] (0.703,0.761) -- (0.285,0.996);
\draw[thick] (0.615,0.604) -- (0.197,0.839);
\end{tikzpicture}
  \task \begin{tikzpicture}[x=0.5cm,y=0.5cm,baseline=(current bounding box.center),line join=round,line cap=round]
\draw[fill=teal!12,draw=teal,thick] (0.000,0.000) -- (3.200,0.000) -- (2.000,1.800) -- (0.000,1.800) -- cycle;
\draw[thin] (0.000,0.440) -- (0.440,0.440) -- (0.440,0.000);
\draw[thin] (0.440,1.800) -- (0.440,1.360) -- (0.000,1.360);
\end{tikzpicture}
\end{tasks}\recorda{\textbf{Paral·lelograms} (costats oposats paral·lels): quadrat, rectangle, rombe i romboide. \textbf{Trapezis} (només dos costats paral·lels): rectangle, isòsceles i escalé. \textbf{Trapezoides}: cap costat paral·lel.}

\ex{1} Tres angles d'un quadrilàter mesuren el que s'indica. Calcula el quart:\begin{tasks}(3)
  \task $90^\circ$, $90^\circ$, $100^\circ$
  \task $75^\circ$, $105^\circ$, $75^\circ$
  \task $120^\circ$, $60^\circ$, $95^\circ$
\end{tasks}\recorda{Els quatre angles d'un quadrilàter sumen $360^\circ$.}

\ex{1} Quins paral·lelograms tenen…\begin{tasks}(2)
  \task les diagonals iguals?
  \task les diagonals perpendiculars?
  \task els quatre angles rectes?
  \task els quatre costats iguals?
\end{tasks}

\ex{2} Calcula els angles que falten:\begin{tasks}(1)
  \task Un rombe amb un angle de $54^\circ$.
  \task Un romboide amb un angle de $112^\circ$.
  \task Un trapezi rectangle amb un angle agut de $63^\circ$.
\end{tasks}

\ex{2} Digues si és cert o fals:\begin{tasks}(1)
  \task Tot quadrat és un rectangle.
  \task Tot rombe és un quadrat.
  \task Tot rectangle és un paral·lelogram.
  \task Un trapezi pot tindre tres angles rectes.
\end{tasks}

\ex{3} Respon:\begin{tasks}(1)
  \task Els angles d'un quadrilàter mesuren $x$, $2x$, $3x$ i $4x$. Quant mesura cadascun?
  \task En un romboide, un angle és el triple del seu consecutiu. Quant mesuren els quatre angles?
\end{tasks}


% =====================================================
\seccio{9}{Angles dels polígons. Polígons regulars}

\ex{1} Quant sumen els angles interiors de cada polígon?\begin{tasks}(5)
  \task pentàgon
  \task hexàgon
  \task octàgon
  \task decàgon
  \task dodecàgon
\end{tasks}\recorda{Els angles d'un polígon de $n$ costats sumen $(n-2)\cdot180^\circ$. En un polígon regular, cada angle interior fa $\dfrac{(n-2)\cdot180^\circ}{n}$ i l'angle central fa $\dfrac{360^\circ}{n}$.}

\ex{1} Calcula l'angle interior i l'angle central d'aquests polígons regulars:\begin{tasks}(3)
  \task triangle equilàter
  \task quadrat
  \task pentàgon
  \task hexàgon
  \task octàgon
  \task decàgon
\end{tasks}

\ex{2} L'angle central d'un polígon regular mesura $24^\circ$. Quants costats té? Quant mesura cada angle interior?

\ex{2} Cada angle interior d'un polígon regular mesura $160^\circ$. Quants costats té?

\ex{2} Els angles interiors d'un polígon sumen $2340^\circ$. Quants costats té?

\ex{3} Pot haver-hi un polígon regular amb angles interiors de $100^\circ$? I de $150^\circ$?

\ex{3} Amb quins d'aquests polígons regulars es pot cobrir el terra sense deixar buits: triangles, quadrats, pentàgons, hexàgons o octàgons? I combinant octàgons i quadrats?

\ex{3} Unim el centre d'un hexàgon regular amb els sis vèrtexs. Quina mena de triangles s'obtenen? Raona-ho amb els angles.

\vfill
\begin{center}\color{muted}\small\faCheckCircle[regular]\ \ Ànim! Si et costa algun exercici, repassa els apunts interactius de la unitat.\end{center}

\end{document}
